% Options for packages loaded elsewhere
\PassOptionsToPackage{unicode}{hyperref}
\PassOptionsToPackage{hyphens}{url}
\documentclass[
]{article}
\usepackage{xcolor}
\usepackage[margin=1in]{geometry}
\usepackage{amsmath,amssymb}
\setcounter{secnumdepth}{-\maxdimen} % remove section numbering
\usepackage{iftex}
\ifPDFTeX
  \usepackage[T1]{fontenc}
  \usepackage[utf8]{inputenc}
  \usepackage{textcomp} % provide euro and other symbols
\else % if luatex or xetex
  \usepackage{unicode-math} % this also loads fontspec
  \defaultfontfeatures{Scale=MatchLowercase}
  \defaultfontfeatures[\rmfamily]{Ligatures=TeX,Scale=1}
\fi
\usepackage{lmodern}
\ifPDFTeX\else
  % xetex/luatex font selection
\fi
% Use upquote if available, for straight quotes in verbatim environments
\IfFileExists{upquote.sty}{\usepackage{upquote}}{}
\IfFileExists{microtype.sty}{% use microtype if available
  \usepackage[]{microtype}
  \UseMicrotypeSet[protrusion]{basicmath} % disable protrusion for tt fonts
}{}
\makeatletter
\@ifundefined{KOMAClassName}{% if non-KOMA class
  \IfFileExists{parskip.sty}{%
    \usepackage{parskip}
  }{% else
    \setlength{\parindent}{0pt}
    \setlength{\parskip}{6pt plus 2pt minus 1pt}}
}{% if KOMA class
  \KOMAoptions{parskip=half}}
\makeatother
\usepackage{graphicx}
\makeatletter
\newsavebox\pandoc@box
\newcommand*\pandocbounded[1]{% scales image to fit in text height/width
  \sbox\pandoc@box{#1}%
  \Gscale@div\@tempa{\textheight}{\dimexpr\ht\pandoc@box+\dp\pandoc@box\relax}%
  \Gscale@div\@tempb{\linewidth}{\wd\pandoc@box}%
  \ifdim\@tempb\p@<\@tempa\p@\let\@tempa\@tempb\fi% select the smaller of both
  \ifdim\@tempa\p@<\p@\scalebox{\@tempa}{\usebox\pandoc@box}%
  \else\usebox{\pandoc@box}%
  \fi%
}
% Set default figure placement to htbp
\def\fps@figure{htbp}
\makeatother
\setlength{\emergencystretch}{3em} % prevent overfull lines
\providecommand{\tightlist}{%
  \setlength{\itemsep}{0pt}\setlength{\parskip}{0pt}}
\usepackage{graphicx}
\usepackage{float}
\usepackage{bookmark}
\IfFileExists{xurl.sty}{\usepackage{xurl}}{} % add URL line breaks if available
\urlstyle{same}
\hypersetup{
  hidelinks,
  pdfcreator={LaTeX via pandoc}}

\author{}
\date{\vspace{-2.5em}}

\begin{document}

\section{Question}\label{question}

Un experimento consiste en medir el alcance horizontal de un proyectil
en función del ángulo con el que se lanza (respecto a la horizontal). En
la gráfica se registran los resultados de 91 lanzamientos realizados con
la misma velocidad inicial.

\includegraphics[width=12cm,height=\textheight,keepaspectratio]{dispersion_proyectil.png}

El comportamiento del alcance respecto al ángulo es

\subsection{Answerlist}\label{answerlist}

\begin{itemize}
\tightlist
\item
  de tipo lineal y con mayor dispersion cuanto mayor sea el angulo de
  lanzamiento.
\item
  de tipo lineal y con mayor dispersion cuanto mayor sea el alcance del
  proyectil.
\item
  parabolico y con mayor dispersion cuanto mayor sea el angulo de
  lanzamiento.
\item
  parabolico y con mayor dispersion cuanto mayor sea el alcance del
  proyectil.
\end{itemize}

\section{Solution}\label{solution}

Para analizar el comportamiento del alcance respecto al ángulo, debemos
observar dos aspectos clave en la gráfica de dispersión:

\textbf{1. Tipo de relacion (de tipo lineal vs parabolico):}

La relacion entre el ángulo y el alcance sigue la ecuación del
movimiento de proyectiles:

\[R = \frac{v_0^2 \sin(2\theta)}{g}\]

Esta es una función senoidal, por lo tanto \textbf{parabolico}. En la
gráfica se observa claramente una forma parabólica con un máximo
aproximadamente en \(\theta \approx 0.78\) radianes (\(\approx 45°\)).

\textbf{2. Patrón de dispersión:}

Observando la gráfica:

\begin{itemize}
\tightlist
\item
  En ángulos bajos (cercanos a 0 radianes): los puntos están más
  concentrados
\item
  En ángulos altos (cercanos a 1.5 radianes): los puntos muestran mayor
  variabilidad
\end{itemize}

Esto indica que \textbf{la dispersión es proporcional a el angulo de
lanzamiento}, no a el alcance del proyectil. Físicamente, esto puede
explicarse porque los errores en la medición del ángulo se amplifican en
configuraciones de lanzamiento con mayor inclinación.

\textbf{Conclusión:}

El comportamiento es \textbf{parabolico} (forma parabólica) y
\textbf{con mayor dispersion cuanto mayor sea el angulo de lanzamiento}
(mayor variabilidad en ángulos altos).

\subsection{Answerlist}\label{answerlist-1}

\begin{itemize}
\tightlist
\item
  \textbf{Falso}. La relación no es de tipo lineal (es parabolico) y
  aunque la dispersión sí depende de el angulo de lanzamiento, el tipo
  de relación es incorrecto.
\item
  \textbf{Falso}. La relación NO es de tipo lineal y la dispersión NO
  aumenta con el alcance del proyectil sino con el angulo de
  lanzamiento.
\item
  \textbf{Verdadero}. La relación es parabolico (forma parabólica) y la
  dispersión aumenta proporcionalmente a el angulo de lanzamiento.
\item
  \textbf{Falso}. Si bien la relación es parabolico, la dispersión NO
  aumenta con el alcance del proyectil sino con el angulo de
  lanzamiento.
\end{itemize}

\section{Meta-information}\label{meta-information}

exname:
dispersion\_alcance\_proyectil\_aleatorio\_interpretacion\_representacion\_n2\_v1
extype: schoice exsolution: 0010 exshuffle: TRUE exsection:
Estadistica/Graficas de Dispersion

\end{document}
